% ============================ SUMMARY SHEET ============================
\begin{center}
{\Large\bfseries The Geometry of a Bicycle City: Behaviour, Rebalancing,\\ Incentives and Fences in Beijing's Dockless Bike-Share System}\\[6pt]
{\large 2026 CQU Mathematical Contest in Modeling, Problem C}\\[10pt]
\begin{tabular}{ll}
\textbf{Team Control Number:} & \rule{3.2cm}{0.4pt} \\
\textbf{Problem Chosen:} & C \\
\end{tabular}
\end{center}

\vspace{4pt}
\hrule
\vspace{8pt}

\begin{center}{\large\bfseries Summary}\end{center}

Beijing's dockless bike-share system in May 2017 was not short of bicycles. It was
short of bicycles \emph{in the right place at the right time}, and it did not know it.
Using 3.21 million order records from 484,754 bicycles ridden by 349,241 users over
fourteen days, this paper builds one behavioural model and uses it to answer four
operational questions on a common footing.

\textbf{The central finding is a diagnosis of oversupply.} We estimate the fleet at
3.4$\times10^5$ bicycles serving 6.2 rides each per day --- a supply level at which
the daily deficit appears as idle stock rather than unmet demand. At every
resolution we examine, from 50 traffic zones down to 5,263 neighbourhood cells, no
zone ever runs out of bicycles. Every subsequent question is therefore posed as a
question of \emph{cost and placement}, not of shortage.

\textbf{Behaviour (Q1).} Trips are modelled as a marked spatio-temporal point
process: a non-homogeneous Poisson arrival stream, a singly-constrained gravity
destination choice with a time-varying distance-decay parameter, and a
speed-based duration model. The decay parameter is identified by maximum
likelihood and ranges from 2.01 (morning, longer commute trips) to 2.52 (midday,
short errands). On two held-out days the model attains a per-trip log-loss of
3.24--3.52 against 3.28--3.65 for a distance-only benchmark and 7.9--8.2 for an
attraction-only one, and it predicts the top-five destinations correctly 40--48\%
of the time. A latent-class analysis separates five user types whose trip shares
(51.8\% for the multi-purpose majority, 29.6\% for a heavy short-trip minority)
explain the shape of the hourly profile.

\textbf{Rebalancing (Q2).} A time-space network linear program moves bicycles
between 50 zones with a fleet of trucks. Because no zone is ever dry, the binding
constraint is not service but the cost of returning the system to a balanced state:
the optimal plan moves 23,959 bicycles per day (1.12\% of all trips) over 5,441
truck-kilometres, at 3.48 CNY per bicycle relocated. Because the trucks run at
night, the plan is best read as an overnight reset. Rebalancing and fleet size are
substitutes: shrinking the fleet to 80\% of its observed size while adding trucks
lowers total daily cost from 427k to 390k CNY, an interior optimum.

\textbf{Incentives (Q3).} We derive the supply curve of voluntary relocations from
a logit acceptance model in which users trade a cash reward against a detour. The
platform should offer a uniform reward of 1.20 CNY, buying 115,385 relocations
(53.3\% of the peak-hour deficit) for 138k CNY and leaving the rest to trucks --- a
34.9\% saving over moving every bicycle by truck. The limit is physical, not
financial: within a 450 m detour only 71.9\% of the deficit has any qualifying
trip nearby, and reaching 95\% requires a 1,025 m catchment.

\textbf{E-fences (Q4).} Parking demand is the arrival process, whose peak we
characterise by Monte Carlo. A greedy maximal-coverage layout of 626 fences at a
450 m walking radius covers 75.3\% of arrivals; the required capacity is brutally
skewed, with a median of 417 slots and a maximum of 4,827. Electronic fences
without enforced caps merely formalise the bicycle mountain.

\textbf{Competition (Q5).} Splitting the fleet among competing platforms does not
reduce availability in a saturated system --- we verify this down to a 35\% share ---
but it multiplies capital, trucks, land and, most damagingly, fragments the data
needed to run any of the models above.

The paper is explicit about what it cannot identify: the initial inventory of
bicycles, the order sampling rate, and the true demand of users who never ride. Each
is bounded rather than assumed away, and the sensitivity analysis reports how far
the conclusions move when they are perturbed.

\vfill
\noindent\textbf{Keywords:} dockless bike sharing; gravity model; time-space network;
linear programming; discrete choice; maximal coverage; Beijing.

\newpage
